\chapter{最后散射与重子生成}

\section{向外看就是向过去看}

在不同的距离上，我们看到的是不同时期的宇宙。在近处，回溯时间并不很重要。一年前或一百年前的宇宙，与现在的宇宙并没有很大不同。但足够远的地方，我们看到的是更年轻的宇宙。更年轻的宇宙更致密，而更致密的宇宙更热。

在某个足够大的距离上，宇宙的温度大约与太阳表面相当。它在我们看来并非如此，因为光子在来到这里的途中已经被红移了。大约一千的红移因子意味着观测到的温度也降低同样的因子：
\[
T_{\mathrm{obs}} \sim \frac{T_{\mathrm{emit}}}{1+z},
\qquad
1+z \sim 10^3.
\]\footnote{这一展示是对00:02:20--00:02:39处口头红移论证的简明编辑转录。}

那就是宇宙变得不透明的地方。宇宙主要成分是氢，还有少量氦。在几千度的温度以上，氢是电离的。因此，在那个表面之外，存在的不是中性原子，而是离子：主要是自由运动着的质子和电子。那里的光子也更有能量，基本上是X射线。如果一个原子偶然形成，它很快又会被碰撞或光子再次电离。

电离气体是电导体，而导体不会让辐射干净地穿过。电荷响应电场，来回振荡，并重新发射辐射。因此，在那个距离之外，宇宙对我们是不透明的。最后散射面就像太阳的可见表面：不是一堵物质墙，而是介质相当突然地从不透明变为透明的地方。

那里的绝大多数粒子是光子，大约多出\(10^8\)倍。今天也是如此，尽管今天的光子能量如此之低，以至于它们不是当前能量密度的重要部分。一旦温度降到氢的电离温度以下，宇宙就变得清澈，我们就能看穿它。

直接的视线到此为止。下一个问题是，已知的原子、核与粒子物理还能把我们带回到多远。

\section{继续向更早回溯}

就在最后散射之内，除了光子之外，主要粒子是质子和电子，还有一些氦核和少量氘核。实际上没有反质子，也实际上没有正电子。但再往更早回溯，温度最终会变得足够高，使两个光子能够碰撞并产生一个电子和一个正电子：
\[
\gamma + \gamma \leftrightarrow e^- + e^+.
\]

基本的阈值是，光子至少必须提供该粒子对的静止能量：
\[
E \ge 2m_e c^2.
\]
低于该阈值，普通光做不到这一点。高于它，电子--正电子对可以被产生，逆反应也可以发生。处于该温度下的一锅热物质，会包含粒子反应，扮演着化学反应在普通物质中所扮演的角色。

刚好高于阈值时，光子仍然比电子和正电子多。电子的静止质量代价高昂。但随着温度升高到远高于阈值，电子的静止能量变得不如动能重要，密度也变得相当。那时电子、正电子和光子都具有大致相同的密度，并疯狂地相互碰撞。

\begin{figure}[ht]
\centering
\includegraphics[width=0.62\textwidth]{lecture_06_figure_02.png}
\caption{比较电子与质子质量尺度的对产生阈值示意图。}
\lectureframe{00:18:47}
\end{figure}

下一个问题是质子会发生什么。它们已经存在，而且到那时已经非常致密，但还没有足够的热能来大量产生质子--反质子对。原因很简单：
\[
m_p \approx 2000\,m_e,
\qquad
E \ge 2m_p c^2
\]
是一个高得多的阈值。\editorialnote{在该图中，较低质量标签不如上方\(2m_e c^2\)那样清晰，但口头比较明确是在电子与质子阈值之间进行的。}

这已经比恒星内部热得多，而在恒星内部，反粒子并不以任何可观的丰度存在。

但即便这也只是一种临时性的描述。当温度高到足以自由地产生质子和反质子时，它也已经高到足以把质子打碎成夸克。因此，真正炽热的汤更应被描述为电子、正电子、夸克、反夸克、光子、中微子、胶子以及其他种类粒子组成的汤。黑板上一个简单的示意性反应是
\[
e^- + e^+ \to \gamma \to q + \bar q.
\]

\begin{classroomqa}
\lecturetimestamp{00:20:05}
\audiencequestion{是介子还是前夸克？}
\lecturerresponse{介子由夸克组成。一旦温度高到足以打碎质子，它也会打碎介子。}
\end{classroomqa}

到目前为止，推理是逆着时间进行的。现在把它反过来，问一问当这锅炽热而对称的汤冷却下来时会发生什么。

\section{电中性以及一个对称的热宇宙}

在足够高的温度下，宇宙包含密度极高的粒子和反粒子。电子与正电子成对出现，夸克与反夸克成对出现，而自然的图景是每种粒子与其反粒子的丰度几乎相等。

电中性需要稍加小心。今天的宇宙是电中性的。一个有用的助记方法是想象一个闭合而有界的宇宙，比如一个三维球面，而不是我们实际的宇宙。先以一个二维球面作为模型来思考。电场线除了在电荷上之外不能起始或终止。正电荷发出力线；负电荷吸收力线。在一个闭合球面上，这些力线不能简单地逃逸到无穷远。它们会在对侧重新汇聚。因此，作为一个助记，一个闭合而有界的宇宙不能带有净电荷。我们实际的宇宙是否字面意义上闭合并不是重点。重点是，大尺度上的正电荷与负电荷似乎抵消得异常之好。

电中性\emph{并不}意味着质子数与反质子数相同。它意味着总正电荷等于总负电荷。质子、正电子以及所有其他带正电的粒子种类合在一起，与电子以及所有其他带负电的粒子种类合在一起相平衡。

在已知物理学能够被信任的范围内，或许还稍微超出一点，宇宙曾经热到足以大量维持粒子和反粒子。任何能叫得出名字的加速器实验，在早期宇宙的某个阶段都有其对应物。一个近乎对称的开端是自然的猜想：电子与正电子数目相等，夸克与反夸克数目相等，并且，如果把质子和反质子用作重子和反重子的代表，
\[
N_p = N_{\bar p}.
\]

这个猜想之所以有吸引力，是因为物理定律在粒子和反粒子之间几乎是对称的。正电子的行为非常像电子，并非完全一样，但非常接近。正是这种近乎对称成了问题所在。

\section{把电影向前放映}

假设宇宙从一个炽热状态开始，其中粒子和反粒子的丰度相等或近乎相等，然后冷却下来。

当温度足够高时，质子和反质子可以通过高能碰撞产生，也可以湮灭回光子。但最终温度会降到质子阈值以下。那时质子--反质子对仍然可以湮灭成光子，但经过一点膨胀之后，那些光子已经损失了足够的能量，无法再变回质子--反质子对。这就是单行道。同样的事情后来也发生在电子和正电子身上，温度比质子阈值低得多，但仍远高于恒星温度。

由于与微观反应速率相比膨胀很慢，在这一过程发生时宇宙大部分时间都接近热平衡。这就是这里所说的绝热膨胀的含义。在这种情形下，更准确的说法是熵变化非常小。对于一锅炽热的粒子汤，熵本质上就是粒子数，只差量级为一的因子。因此，湮灭掉的质子、反质子、电子和正电子大致被光子所取代。

如果质子数与反质子数完全相等，它们就会直接湮灭殆尽。电子和正电子也会如此。人们将只剩下光子。但如果质子比反质子哪怕只有极微小的过剩，成对的粒子会湮灭，而未配对的过剩部分会留下来。一个具体的例子可以说明这一点：如果有一百个质子和五十个反质子，那五十个反质子会找到五十个质子，与它们湮灭，而五十个质子会留下来。电荷中性随后要求电子相对于正电子也有同样类型的剩余过剩。

这使人们能够从今天的光子数读出早期的不对称。今天光子数大约是质子数或电子数的\(10^9\)到\(10^{10}\)倍：
\[
\frac{N_\gamma}{N_p} \sim 10^9\text{--}10^{10}.
\]
所以在早期炽热的汤中，每一十亿个质子--反质子对大约只有一个多余的质子。正是这种微小的不平衡，构成了原子至今仍然存在的全部原因。

\begin{classroomqa}
\lecturetimestamp{00:42:25}
\audiencequestion{为什么不从统计的角度来考虑这个过剩，把它看作是在冷却过程中以极小概率产生的东西呢？}
\lecturerresponse{要从一个初始对称的状态中得到一个额外的质子，就必须在不产生反质子的情况下产生一个质子。同时还必须产生与之相伴的负电荷，比如说一个电子。已知的标准模型中不包含这样的低能过程。在那些已被检验的能量下，重子数是守恒的，因此普通的冷却过程不会产生净过剩。}
\end{classroomqa}

所以电荷并不是障碍。一个示意性的例子是
\[
\gamma + \gamma \to p + e^-.
\]\footnote{这个展示是对00:44:16--00:46:35处口头例子的编辑转写。其目的是分离出这样一个要点：电荷可以守恒，而重子数被破坏。}
这在整体上会是电中性的。但它违反了另一条经验规则。在被检验的标准模型中，不可能在不产生反质子的情况下产生一个质子。就目前的目的而言，相关的守恒量是重子数，
\[
B=(N_p+N_n)-(N_{\bar p}+N_{\bar n}).
\]\footnote{这个紧凑的公式形式化了00:45:18--00:45:40处给出的口头定义。}
在已知的能量下，该量是守恒的，因此这种不对称不可能在冷却历史中晚期、已被充分检验的那一部分中产生出来。

\section{重子生成}

这就是所谓的重子生成。假设宇宙确实是以等量的重子和反重子开始的。那么物理定律必须允许什么，才能把那个对称状态倾斜成观测到的微小过剩呢？

在黑板上，这被归结为两个符号表达式，
\[
N_p=N_{\bar p},
\qquad
N_p-N_{\bar p},
\]
然后归结为一系列把第一个与第二个联系起来的条件。

\begin{figure}[ht]
\centering
\includegraphics[width=0.72\textwidth]{lecture_06_figure_04.png}
\caption{从一个初始对称状态产生重子不对称的条件。}
\lectureframe{01:22:26}
\end{figure}

\editorialnote{在这个画面中，编号列表中的文字只有一部分可以辨认。下面三项依据的是随附的口头讨论。}

第一个条件是重子数的破坏。如果重子数是一条精确的守恒定律，那么一个以零净重子数开始的宇宙就会一直保持那样。

第二个条件是粒子与反粒子之间存在偏向。仅有重子数破坏还不够。如果自然界允许一个产生质子的过程及其镜像的产生反质子的过程以相等的概率发生，那么平均失衡仍然为零。需要的是两者的速率不同。

普通的统计涨落也远远太小。如果粒子总数约为\(10^{90}\)，包括光子在內，那么量级为\(\sqrt{N}\)的涨落只能给出大约\(10^{45}\)，远不足以解释宇宙中大约\(10^{80}\)个质子。

第三个条件是宇宙必须在某个时刻足够偏离平衡。这里的论证用到了CPT定理。如果一团热气体永远处于平衡之中，那么倒放的影片看起来与正放的影片相同。在这种情形下，粒子与反粒子之间的对称性被那个把电荷共轭、宇称和时间反演结合起来的精确对称性有效地恢复。要得到一个真正被冻结下来的偏向，实际的历史在微观时间尺度上正放和倒放时必须看起来不同。足够早的时候，宇宙的膨胀足够快，可以做到这一点。

所以这三个条件是：
\begin{enumerate}
\item 重子数的破坏；
\item 粒子与反粒子之间对称性的破坏，更确切地说是CP破坏；
\item 一个偏离热平衡的阶段。
\end{enumerate}

\begin{classroomqa}
\lecturetimestamp{01:09:08}
\audiencequestion{把这些与CPT对应起来，条件2真的是CP，而条件3真的是时间方向的问题吗？}
\lecturerresponse{是的。条件2实际上应当表述为CP破坏。条件3是实际历史偏离平衡这一事实，而不是一条单独的守恒定律，也不等同于破坏重子数。}
\end{classroomqa}

有了这三个条件，即使是一个以等量质子和反质子开始的宇宙，也可以被朝一个方向或另一个方向倾斜。仍然不确定的是，自然界是否真的破坏重子数，如果是的话，又是如何破坏的。

\section{已知的是什么，未知的是什么}

这个故事尚未解决的部分是重子数破坏。我们从未在自然界中直接看到过改变
\[
(N_p+N_n)-(N_{\bar p}+N_{\bar n}).
\]
的反应。这可能只意味着相关反应需要比实验室中直接测试过的任何温度都高得多的温度，或者它们在低能下极其不可能发生。但目前这仍然是一个问号。

这就是为什么质子衰变在讨论中如此重要。如果重子数并非严格守恒，那么质子在原则上应该是不稳定的。但它却极其稳定。这里引用的下限约为
\[
\tau_p \gtrsim 10^{34}\ \text{years}.
\]
这来自使用极大量水、配备光探测器并深埋地下的大型实验。它并未否定重子数破坏。它只意味着任何此类破坏在低能下都必须受到极其强烈的压制。

还有一个已知物理学的评论也进入了讨论。在\(100\)到\(250\,\mathrm{GeV}\)附近，弱相互作用变得与电磁相互作用更加相似。这在详细模型中可能很重要，但在当前论证中并不是戏剧性的一步。

暗物质并没有在很大程度上改变基本逻辑。在大多数理论中，它与普通物质的相互作用如此微弱，以至于对于这个问题而言，它基本上只是一个旁观者。

\begin{classroomqa}
\lecturetimestamp{01:24:49}
\audiencequestion{引力如何影响这一切？}
\lecturerresponse{主要通过两种方式。第一，没有宇宙膨胀就不会有非平衡阶段。第二，黑洞蒸发给出了一个强有力的理由，使人怀疑重子数不可能永远是一条精确的守恒定律。除此之外，粒子之间的普通引力相互作用太弱，在这里无关紧要。}
\end{classroomqa}

接着有一个关于电荷的后续问题。同样的黑洞论证是否也会威胁电荷守恒？答案是否定的。带电黑洞携带着向外延伸的电通量线，而这些线记录着电荷。这样的黑洞不会一直蒸发殆尽；相反，它会趋近于一个极端残余。因此，黑洞论证威胁的是重子数，而不是电荷。

另一个问题问，极其致密的早期夸克汤是否应被简单地视为黑洞。答案是否定的。要点并不是一个静态的致密团块。宇宙当时正在膨胀。黑洞故事是一个关于重子数命运的思维实验，而不是把早期宇宙字面描述为一个被囚禁的区域。

\begin{classroomqa}
\lecturetimestamp{01:31:14}
\audiencequestion{人们会知道如何修改标准模型以包含此类反应吗？}
\lecturerresponse{麻烦恰恰相反。很容易加入过多的重子数破坏，使质子衰变得太快。真正的问题是在仍然允许人们想要的早期宇宙过程的同时，让质子保持像现在这样稳定。}
\end{classroomqa}

\section{熵与晚期中子问题}

现在回到记账。在绝热膨胀的宇宙中，熵近似守恒。对于基本粒子的热汤，熵本质上就是粒子数，最多相差数量级为一的因子。这就是为什么讨论不需要凭空在晚期产生大量光子。湮灭主要是把汤的内容重新组合成辐射。

这个问题的旧版本把故事讲反了。它不是问为什么相对于光子只有如此少的质子存活下来，而是反过来问：为什么相对于质子有如此多的光子？但更尖锐的问题并不是为什么宇宙在后期制造了如此多的熵。而是为什么粒子与反粒子之间的不对称从一开始就如此微小：大约每十亿对质子--反质子对才多出一个质子。

自由中子确实会衰变，但该过程\emph{并不}破坏重子数：
\[
n \to p + e^- + \bar\nu.
\]
中子携带与质子相同的重子数。

对于自由中子，可用的衰变能量非常小，因为中子质量只略高于末态质量。这也是为什么按基本粒子的标准该衰变如此缓慢的部分原因。当中子束缚在原子核中时，结合能可以完全消除那一点可用能量，于是束缚中子就可以稳定。

\begin{classroomqa}
\lecturetimestamp{01:43:03}
\audiencequestion{中子是在从宇宙中简单地消失，还是也会被产生？}
\lecturerresponse{第一个答案是自由中子通过衰变而损失。但这并不是故事的全部。存在真实的中子产生来源，包括恒星中的核反应和中子星的形成。因此中子总数并不是简单地单向减少。}
\end{classroomqa}
